\section{Related Work}
\label{sec:related}

\paragraph{Finding cryptography in software.} Static analysis of
cryptographic code began with misuse: CryptoLint found that most Android apps
using cryptographic APIs misused them~\cite{egele2013}, and CrySL and
CryptoGuard check Java code against usage rules at scale~\cite{kruger2021crysl,
rahaman2019}, motivated in part by how hard developers find these
APIs~\cite{nadi2016}. These tools ask whether cryptography is used
correctly; an inventory asks what is used and where, including correct uses of
algorithms a quantum computer breaks. The open-source inventory tools we
compare in \S\ref{sec:rq1} (CBOMkit~\cite{cbomkit},
cryptoscan-pqc~\cite{cryptoscanpqc}, pqc-migration-audit~\cite{pqcaudit} and
Semgrep's rules~\cite{semgrep}) each cover one or two layers.
Olewinski et al.\ survey open-source cryptographic asset discovery tools and
find the landscape fragmented across categories~\cite{olewinski2026}; our
corpus measures that fragmentation as recall per layer. Other work covers
single layers in depth: Kestrel detects ML-KEM and ML-DSA implementations in
stripped binaries from their constant tables~\cite{junaid2026}, and Cho et al.\
flag quantum-vulnerable TLS from captured traffic~\cite{cho2026}. \tool{}'s
contribution is the combination: five layers normalized to one finding type
and reconciled, so that a gap between layers becomes a finding of its own.
For dependencies we follow the lesson of vulnerability analysis that a
library's presence is not its use~\cite{ponta2018}, and report dependency
findings as capabilities.

\paragraph{Post-quantum readiness of open source.} Yuan et al.\ measure
post-quantum support across eight package ecosystems and find it
rare~\cite{yuan2026}. RQ3 asks the converse question for the most critical
projects: how much quantum-vulnerable cryptography they carry, in which
layer, and how exposed it is. Library-level surveys track which libraries
offer the NIST algorithms~\cite{ahmed2025}, which is what \tool{}'s
dependency layer encodes.

\paragraph{Measurements of post-quantum handshakes.} Wickramasinghe et
al.~\cite{wickramasinghe2026} measure post-quantum TLS longitudinally across a
million domains from eleven vantage points, find that deployment centers on a
single hybrid construction driven by managed infrastructure providers, and see
no meaningful latency cost. Our results agree on all three. Dubey and
Varshney report hybrid key exchange on about half of a cross-sector domain list
and no post-quantum certificates~\cite{dubey2026}. Loizou and Ghadafi measure
UK organizations' HTTPS and mail endpoints and find, as we do, that the
provider predicts deployment better than the organization's
sector~\cite{loizou2026}. Haque et al.\ extend the measurement beyond HTTPS to
directory, mail and remote-access protocols, where support is lower,
mostly far lower~\cite{haque2026}. We differ in population and scope: the complete
registry of federal domains that EO~14412 and M-26-15 govern, broken down by
agency as well as provider, and paired with certificate and code-level
measurements. Singanamalla et al.\ measured the long tail of government adoption of
HTTPS itself~\cite{singanamalla2020}, and
certificate measurement at Internet scale goes back to
Durumeric et al.~\cite{durumeric2013https}.

\paragraph{Cost of post-quantum handshakes.} Paquin et al.\ benchmark post-quantum
key exchange and signatures in TLS under emulated network
conditions~\cite{paquin2020}, Sikeridis et al.\ study post-quantum
authentication in TLS~1.3~\cite{sikeridis2020}, and Kampanakis and
Childs-Klein measure the effect on time to last byte of real
connections~\cite{kampanakis2024}. The sizes of the NIST schemes follow from
their specifications~\cite{fips203,fips204,fips205}. RQ4 repeats the
measurement for the standardized hybrid group in a production TLS stack, the
one our probe is checked against, and cross-checks every size against OpenSSL.

\paragraph{Migration and crypto-agility.} A community workshop report frames
migration and crypto-agility as research problems~\cite{ott2019}. NIST's NCCoE migration project~\cite{nccoe} and CISA's
strategy for automated discovery and inventory tools~\cite{cisaacdi} make
discovery the starting point for federal agencies, which is the task \tool{}
automates.

\paragraph{Measurement ethics.} Our scanning practice follows the guidance
of the ZMap authors~\cite{zmap} and the Menlo Report~\cite{menlo}.
